%% ============================================================
%%  INTRODUCTION (chapitre non numéroté)
%% ============================================================
\chapter*{Introduction}
\addcontentsline{toc}{chapter}{Introduction}
\markboth{INTRODUCTION}{INTRODUCTION}

À une époque où la technologie reconfigure chaque aspect de l'apprentissage
et du travail, l'intelligence artificielle générative s'impose comme l'une des
forces les plus transformatrices de notre temps. Ce qui a commencé par des
systèmes expérimentaux capables de produire des textes ou des images simples a
évolué vers des outils sophistiqués capables de créer du contenu, de raisonner
sur des problèmes complexes et d'interagir de manière dynamique avec des systèmes
du monde réel. Pour un étudiant en master professionnel de cybersécurité, explorer
l'IA générative ne relève pas d'une simple curiosité académique. Il s'agit d'une
opportunité concrète pour répondre à des enjeux réels en technologie éducative,
en particulier dans un contexte d'accélération de la transformation numérique des
établissements scolaires.

Le parcours de l'IA générative remonte plus loin qu'on ne le pense souvent. Ses
racines se trouvent dans les premières expériences computationnelles du milieu du
\textsc{xx}\textsuperscript{e} siècle, lorsque les chercheurs ont commencé à imaginer
des machines capables d'imiter la créativité humaine. Dès les années 1960, des
chatbots rudimentaires comme ELIZA ont mis en évidence le potentiel des machines
à générer des réponses conversationnelles, même si elles s'appuyaient sur une
simple correspondance de motifs plutôt que sur une véritable compréhension.

La véritable rupture est intervenue dans les années 2010 avec l'essor des techniques
d'apprentissage profond. En 2014, Ian Goodfellow et ses collègues ont introduit les
réseaux antagonistes génératifs (GANs), un cadre dans lequel deux réseaux neuronaux
entrent en compétition pour produire des données de plus en plus réalistes,
qu'il s'agisse d'images, d'audio ou de texte~\cite{goodfellow2014}. Cette innovation
a ouvert la voie à la production de contenu synthétique de haute qualité et posé les
principes fondateurs des modèles ultérieurs.

Le rythme s'est ensuite accéléré avec l'introduction de l'architecture Transformer
en 2017. L'article \og{}Attention Is All You Need\fg{} d'Ashish Vaswani et de son
équipe a révolutionné le traitement automatique du langage en remplaçant les réseaux
récurrents par des mécanismes d'auto-attention capables de traiter les séquences en
parallèle et de capturer beaucoup plus efficacement les dépendances à longue
portée~\cite{vaswani2017}. Cette architecture est devenue la colonne vertébrale de
la série GPT développée par OpenAI. GPT-1 est apparu en 2018, suivi de GPT-2 en 2019
puis du modèle beaucoup plus volumineux GPT-3 en 2020. Chaque itération a augmenté le
nombre de paramètres et l'ampleur des données d'entraînement, produisant des sorties
d'apparence de plus en plus humaine.

L'explosion auprès du grand public a eu lieu en novembre 2022 avec ChatGPT, une
interface conversationnelle construite sur GPT-3.5 qui a fait entrer l'IA générative
dans des millions de foyers et de salles de classe du jour au lendemain. Tout à coup,
n'importe qui pouvait générer des essais, du code, des plans de cours ou des histoires
créatives à partir d'une simple demande en langage naturel. En 2023 et 2024, le domaine
a dépassé le cadre des simples chatbots mono-modèle pour s'orienter vers des systèmes
multimodaux et, surtout, vers l'IA agentique : des systèmes capables de planifier,
d'utiliser des outils et d'interagir avec des environnements externes. Des modèles
comme GPT-4o, Claude ou Grok ont intégré des capacités de vision, de voix et de
raisonnement.

Toutefois, l'augmentation d'échelle seule s'est révélée insuffisante. Les chercheurs
ont observé des limitations récurrentes : hallucinations factuelles, absence d'accès
à l'information à jour et difficultés d'interaction fiable avec les systèmes logiciels
existants. Ces défis ont conduit à l'émergence de techniques complémentaires, comme la
Retrieval-Augmented Generation (RAG) pour ancrer les réponses dans des données réelles
et, plus récemment, le Model Context Protocol (MCP). Introduit par Anthropic en novembre
2024, MCP établit une norme ouverte permettant aux modèles d'IA de découvrir et de
contacter de manière sécurisée des outils, des bases de données et des services
externes~\cite{anthropic2024}.

%% ── Figure : frise chronologique ───────────────────────────
\begin{figure}[htbp]
\centering
\resizebox{\linewidth}{!}{%
\begin{tikzpicture}[
    timeline/.style={line width=3pt, draw=gray!60},
    milestone/.style={circle, draw=gray!70, thick, minimum size=11mm,
      inner sep=2pt, fill=white, font=\small\bfseries},
    label/.style={align=center, font=\footnotesize}
]
\draw[timeline] (0,0) -- (16.5,0);
\node[milestone] (e1) at (0.8,0)  {1960s};
\node[label, above=12pt] at (e1)  {ELIZA\\(premier chatbot)};
\node[milestone] (e2) at (3.8,0)  {2014};
\node[label, above=12pt] at (e2)  {GANs};
\node[milestone] (e3) at (6.5,0)  {2017};
\node[label, above=12pt] at (e3)  {Transformers};
\node[milestone] (e4) at (9.2,0)  {2020};
\node[label, above=12pt] at (e4)  {GPT-3};
\node[milestone] (e5) at (11.8,0) {2022};
\node[label, above=12pt] at (e5)  {ChatGPT};
\node[milestone] (e6) at (15.2,0) {2024};
\node[label, above=12pt] at (e6)  {MCP \&\\IA agentique};
\draw[-{Stealth[length=3.5mm]}, gray!70, thick] (1.9,0)  -- (3.2,0);
\draw[-{Stealth[length=3.5mm]}, gray!70, thick] (4.9,0)  -- (6.0,0);
\draw[-{Stealth[length=3.5mm]}, gray!70, thick] (7.6,0)  -- (8.7,0);
\draw[-{Stealth[length=3.5mm]}, gray!70, thick] (10.3,0) -- (11.3,0);
\draw[-{Stealth[length=3.5mm]}, gray!70, thick] (12.9,0) -- (14.7,0);
\end{tikzpicture}%
}%
\caption{Évolution de l'IA générative : jalons clés}
\label{fig:timeline}
\end{figure}

%% ────────────────────────────────────────────────────────────
L'entreprise Intellect Education a développé SprintUp, une plateforme innovante qui
fonctionne comme un \og{}Shopify de l'éducation\fg{}. Concrètement, SprintUp permet
aux écoles et aux institutions de générer en quelques minutes des campus numériques
entièrement fonctionnels : portails étudiants, tableaux de bord enseignants, gestion
de curriculum, classes en direct et services backend exposant des API REST.

SprintUp intègre déjà une fonctionnalité de génération de syllabus fondée sur un
appel direct à un modèle de langage unique (LLM). Cependant, cette approche
mono-modèle présente des limitations significatives en termes de coût, de rapidité
et de qualité du contenu produit.

\subsection*{Problématique}

L'approche actuelle de génération de curriculum sur SprintUp repose
sur un \textbf{appel unique} à un \gls{llm} pour produire
l'ensemble du contenu. Ce fonctionnement souffre de trois
limitations majeures qui touchent toutes à la \textbf{qualité} du
contenu produit.

Le premier problème concerne la \textbf{redondance}. Un appel
mono-modèle produit du contenu nouveau à chaque exécution, sans
mémoire de ce qui a déjà été produit auparavant. Sur un parcours
pédagogique constitué au fil de plusieurs sessions, plusieurs
activités finissent par traiter le même concept sous des angles à
peine différents, ce qui dilue l'ancrage des objectifs et
alourdit inutilement la charge cognitive de l'apprenant. Le
deuxième porte sur l'\textbf{ancrage factuel} : en l'absence de
mécanisme de contrôle, le contenu produit varie d'un appel à
l'autre et aucun garde-fou n'empêche les hallucinations, les
sections manquantes ou les contenus trop superficiels. Le
troisième est l'\textbf{absence de séparation des privilèges} :
un même prompt donne accès à toutes les opérations, ce qui rend
tout usage spécialisé (révision par un tuteur, accompagnement
d'un apprenant) difficile à isoler et à sécuriser.

\medskip
\noindent
La question qui structure ce travail est donc la suivante :
\emph{comment concevoir une famille d'agents spécialisés,
intégrés de manière standardisée à une base de connaissances
pédagogique partagée, qui produit un contenu sans redondance par
rapport à l'existant et dont chaque agent dispose du périmètre
d'accès strictement nécessaire à son rôle ?}

L'objectif central de ce projet est de remplacer l'approche
mono-modèle par une \textbf{famille d'agents spécialisés}, reliés
entre eux par un orchestrateur dédié et reliés à la base de
données pédagogique par un serveur Model Context Protocol commun.
Les agents sont au nombre de cinq dans cette première version :
un \textbf{générateur d'activités sans \gls{mcp}} qui sert de
baseline, un \textbf{générateur d'activités avec \gls{mcp} et
indexation} qui consulte le contenu existant avant d'en produire,
un \textbf{générateur de syllabus} qui compose ce contenu en
parcours pédagogiques cohérents, un \textbf{tuteur} qui révise
les productions d'apprenants et un \textbf{assistant élève} qui
répond aux questions ancré dans le syllabus suivi. Un sixième
agent, l'\textbf{assistant administrateur}, est mentionné en
perspective.

Plus précisément, ce travail poursuit trois objectifs. Le
premier est la conception d'un \textbf{mécanisme d'indexation}
des activités existantes, sur lequel les agents générateurs
s'appuient pour éviter de produire du contenu redondant. Le
deuxième est la conception d'un \textbf{serveur \gls{mcp}} qui
expose ces activités et les opérations associées sous la forme
d'un catalogue d'outils typés, indépendamment de la nature des
clients qui les consomment. Le troisième est la conception du
\textbf{catalogue d'agents} décrivant les rôles retenus pour
cette première version, avec pour chacun une liste blanche
d'outils \gls{mcp} qui matérialise le principe du moindre
privilège. La \textbf{stratégie d'orchestration} qui articule la
coexistence de ces agents, ainsi que le dispositif d'observabilité
et l'analyse de sécurité associés, font l'objet du
\textbf{chapitre~3}, qui prolonge directement la présente conception.

Le chapitre~1 présente les cadres d'orchestration d'agents, les
techniques d'augmentation par recherche et le Model Context
Protocol, qui constituent les fondations sur lesquelles s'appuie la
conception du système. Le chapitre~2 détaille cette conception —
indexation des activités pédagogiques, protocole \gls{mcp} appliqué
au système et catalogue des agents retenus pour cette première
version — puis en propose une lecture comparative. Il s'achève en
annonçant le \textbf{chapitre~3}, consacré à la stratégie
d'orchestration, à l'observabilité et à la sécurité applicative,
qui constitue la suite directe du présent travail.

En résumé, ce projet se situe à l'intersection entre l'\gls{ia}
générative de pointe et des besoins éducatifs concrets. En
construisant une famille d'agents spécialisés autour d'une base
de connaissances pédagogique partagée et d'un mécanisme
d'indexation commun, nous cherchons à démontrer qu'une
architecture orientée \textbf{qualité} produit un contenu plus
cohérent et plus ancré que l'approche mono-modèle existante,
tout en respectant les exigences de cybersécurité d'un
environnement éducatif.
\newpage
